\documentclass[11pt,a4paper]{article}
\usepackage[utf8]{inputenc}
\usepackage[T1]{fontenc}
\usepackage[spanish,es-noshorthands]{babel}
\usepackage[margin=2.5cm]{geometry}
\usepackage{amsmath,amssymb}
\usepackage{longtable}
\usepackage{booktabs}
\usepackage{array}
\usepackage{enumitem}
\usepackage{hyperref}
\hypersetup{colorlinks=true, linkcolor=blue, citecolor=blue, urlcolor=blue}

\title{Cronograma clase por clase \\ \large Probabilidad y Estadística (5765)}
\author{Universidad Nacional del Sur}
\date{}

\begin{document}
\maketitle
\vspace{-1.5em}

\begin{center}
\textit{Cronograma tentativo, proyectado suponiendo dictado todos los martes y jueves desde el 18 de agosto de 2026, sin interrupciones. Sujeto a los ajustes que disponga la cátedra.}
\end{center}
\vspace{1em}

\section*{Feriados relevantes del período (Argentina 2026, verificados)}

Se revisó el calendario oficial de feriados nacionales 2026 y el feriado local de Bahía Blanca (aniversario de fundación, 11 de abril, fuera del período de cursada) contra todas las fechas de martes y jueves entre el 18/8 y el 3/12/2026. Ninguna clase programada coincide con un feriado. Para referencia, los feriados nacionales del período y su día de la semana:

\begin{center}
\begin{tabular}{@{}lll@{}}
\toprule
\textbf{Fecha} & \textbf{Feriado} & \textbf{Día / observación} \\
\midrule
Lunes 17/8 & Paso a la Inmortalidad del Gral. San Martín & Lunes --- no coincide con clase \\
Lunes 12/10 & Día del Respeto a la Diversidad Cultural & Lunes --- no coincide con clase \\
Viernes 20/11 & Día de la Soberanía Nacional (fecha original) & Viernes --- no coincide con clase \\
Lunes 23/11 & Día de la Soberanía Nacional (trasladado) & Lunes --- no coincide con clase \\
Lunes 7/12 & Puente turístico (no laborable) & Lunes --- no coincide con clase \\
\textbf{Martes 8/12} & \textbf{Inmaculada Concepción (feriado fijo)} & \textbf{Martes --- ¡atención!} \\
\bottomrule
\end{tabular}
\end{center}

\vspace{0.5em}
\begin{center}
\parbox{0.9\textwidth}{\small \textbf{Nota:} el martes 8/12 es feriado nacional fijo. La última clase de teoría (Clase 32) cae el jueves 3/12, antes del feriado, por lo que no hay conflicto con el dictado de clases. Sí conviene tenerlo en cuenta al fijar la fecha del 2\textsuperscript{do} parcial o del recuperatorio si se piensan programar para la semana del 7 al 11 de diciembre: \textbf{evitar el martes 8/12}.}
\end{center}

\newpage
\section*{Calendario de clases}

\begin{longtable}{@{}c l l l@{}}
\toprule
\textbf{Clase} & \textbf{Fecha} & \textbf{Día} & \textbf{Unidad} \\
\midrule
\endfirsthead
\toprule
\textbf{Clase} & \textbf{Fecha} & \textbf{Día} & \textbf{Unidad} \\
\midrule
\endhead
\midrule \multicolumn{4}{r}{\textit{Continúa en la página siguiente}} \\
\endfoot
\bottomrule
\endlastfoot

Clase 0 & 18/ago & Martes & Presentación de la materia \\
Clase 1 & 18/ago & Martes & Unidad 1: Espacios de probabilidad \\
Clase 2 & 20/ago & Jueves & Unidad 1: Espacios de probabilidad \\
Clase 3 & 25/ago & Martes & Unidad 1: Espacios de probabilidad \\
Clase 4 & 27/ago & Jueves & Unidad 1: Espacios de probabilidad \\
Clase 5 & 1/sep & Martes & Unidad 2: Variables aleatorias \\
Clase 6 & 3/sep & Jueves & Unidad 2: Variables aleatorias \\
Clase 7 & 8/sep & Martes & Unidad 2: Variables aleatorias \\
Clase 8 & 10/sep & Jueves & Unidad 2: Variables aleatorias \\
Clase 9 & 15/sep & Martes & Unidad 2: Variables aleatorias \\
Clase 10 & 17/sep & Jueves & Unidad 3: Vectores aleatorios \\
Clase 11 & 22/sep & Martes & Unidad 3: Vectores aleatorios \\
Clase 12 & 24/sep & Jueves & Unidad 3: Vectores aleatorios \\
Clase 13 & 29/sep & Martes & Unidad 3: Vectores aleatorios \\
Clase 14 & 1/oct & Jueves & Unidad 4: Función generadora de momentos \\
Clase 15 & 6/oct & Martes & Unidad 4: Función generadora de momentos \\
Clase 16 & 8/oct & Jueves & Unidad 4: Función generadora de momentos \\
Clase 17 & 13/oct & Martes & Unidad 4: Función generadora de momentos \\
Clase 18 & 15/oct & Jueves & Unidad 4: Función generadora de momentos \\
\multicolumn{4}{l}{\rule{0pt}{2.2ex}\textit{1\textsuperscript{er} parcial (tentativo, fecha a confirmar) --- última semana de octubre}} \\[0.3em]
Clase 19 & 20/oct & Martes & Unidad 5: Muestras aleatorias y distribuciones \\
Clase 20 & 22/oct & Jueves & Unidad 5: Muestras aleatorias y distribuciones \\
Clase 21 & 27/oct & Martes & Unidad 6: Estimación de parámetros \\
Clase 22 & 29/oct & Jueves & Unidad 6: Estimación de parámetros \\
Clase 23 & 3/nov & Martes & Unidad 6: Estimación de parámetros \\
Clase 24 & 5/nov & Jueves & Unidad 6: Estimación de parámetros \\
Clase 25 & 10/nov & Martes & Unidad 7: Prueba de hipótesis \\
Clase 26 & 12/nov & Jueves & Unidad 7: Prueba de hipótesis \\
Clase 27 & 17/nov & Martes & Unidad 7: Prueba de hipótesis \\
Clase 28 & 19/nov & Jueves & Unidad 7: Prueba de hipótesis \\
Clase 29 & 24/nov & Martes & Unidad 8: Modelos lineales \\
Clase 30 & 26/nov & Jueves & Unidad 8: Modelos lineales \\
Clase 31 & 1/dic & Martes & Unidad 8: Modelos lineales \\
Clase 32 & 3/dic & Jueves & Unidad 8: Modelos lineales \\
\multicolumn{4}{l}{\rule{0pt}{2.2ex}\textit{2\textsuperscript{do} parcial y recuperatorio (tentativo, fecha a confirmar) --- diciembre, evitando el 8/12}} \\[0.3em]

\end{longtable}

\vspace{1em}
{\small \textbf{Aclaraciones:} el martes 18/8 concentra la Clase 0 (presentación) y la Clase 1 (primer tema de la Unidad 1) en el mismo día. A partir de ahí, cada martes y jueves corresponde a una clase de teoría; las prácticas se dictan los jueves antes de la teoría (ver Clase 0 para más detalle), y hay un horario adicional de práctica todavía a confirmar. Las fechas exactas de los parciales y el recuperatorio quedan sujetas a lo que confirme la asistente de cátedra; este documento se actualizará cuando estén definidas.}

\end{document}
